\documentclass[12pt]{article}
\usepackage[T1]{fontenc}
\usepackage[utf8]{inputenc}
\usepackage{lmodern}
\usepackage{textcomp}
\usepackage{enumitem}
\usepackage{geometry}
\usepackage{graphicx}
\usepackage{amsmath, amssymb}
\geometry{margin=1in}
\usepackage{calc}
\usepackage{xcolor}
% --- latex-dual-layer macros (auto-generated) ---
\newlength{\dlboxwidth}
\newlength{\dlboxheight}
\newlength{\dlboxdepth}
\newcommand{\duallayerbox}[2]{%
  \begingroup
  \settowidth{\dlboxwidth}{\strut #1}%
  \settoheight{\dlboxheight}{\strut #1}%
  \settodepth{\dlboxdepth}{\strut #1}%
  \ifdim\dlboxwidth=0pt
    \settowidth{\dlboxwidth}{#2}%
  \fi
  \raisebox{0pt}[\dlboxheight][\dlboxdepth]{%
    \makebox[\dlboxwidth][l]{\resizebox{\dlboxwidth}{!}{\strut #2}}%
  }%
  \endgroup
}
% --- end latex-dual-layer macros ---
\begin{document}
\begin{center}
Religious Studies - K-12 Level Assessment 1 \
Total Marks: 40 \
Answer all questions.
\end{center}

\section*{Multiple Choice (10 marks)}
\begin{enumerate}[label=\arabic*.]
\item What is a \duallayerbox{mohel}{menorah}?
\begin{enumerate}[label=(\alph*)]
    \item A \duallayerbox{ritual circumciser}{fisherman}
    \item A type of oil lame
    \item A teacher
    \item A covenant
\end{enumerate}
\item Rome claimed that their bishop (pope) was the director successor of which leader?
\begin{enumerate}[label=(\alph*)]
    \item Peter
    \item Paul
    \item Jesus
    \item Matthew
\end{enumerate}
\item When was the first Buddhist temple constructed in Japan?
\begin{enumerate}[label=(\alph*)]
    \item 325 CE
    \item 119 CE
    \item 451 CE
    \item 596 CE
\end{enumerate}
\item Which of the following was part of the role of a deaconess?
\begin{enumerate}[label=(\alph*)]
    \item Ministering to the sick
    \item Preparing women for baptism
    \item Praying for the suffering
    \item All of the above
\end{enumerate}
\item The term xin refers to which of the following?
\begin{enumerate}[label=(\alph*)]
    \item Piety
    \item Non-action
    \item Heart-mind
    \item World
\end{enumerate}
\end{enumerate}

\section*{True / False (10 marks)}
\textit{Answer True or False}
\begin{enumerate}[label=\arabic*.]
\setcounter{enumi}{5}
\item True or False: In Sikhism, the communal meal offered at the place of worship is called \duallayerbox{Sangat}{Langar}.
\item True or False: Bhakti is often translated as \duallayerbox{yoga}{devotion}, emphasizing physical practices rather than emotional connection to the divine.
\item True or False: The Celestial Masters are an \duallayerbox{anti-Han rebel group}{religious movement} that belongs to the Daoist tradition.
\item True or False: Baisakhi Day is the \duallayerbox{most important}{a significant} festival for Sikhs, celebrating the harvest and the formation of the Khalsa.
\item True or False: In Sikhism, \duallayerbox{Karam}{Hukam} refers to the divine order as an all-embracing higher principle governing life.
\end{enumerate}

\section*{Long Form Response (20 marks)}
\textit{Answer each question with a clear and complete explanation.}
\begin{enumerate}[label=\arabic*.]
\setcounter{enumi}{10}
\item Discuss the significance of the \duallayerbox{Akitu festival}{Zagmuk festival} in ancient Babylonian culture, including
its religious, social, and cultural implications during the New Year celebration.
\item Discuss how the \duallayerbox{Egyptian composition Ludul Bel Nemequi}{Babylonian text Enuma Elish} reflects the dual nature of Marduk
as both wrathful and merciful, providing examples to support your analysis.
\end{enumerate}

\vfill
\noindent\textit{End of Paper}
\end{document}